% Options for packages loaded elsewhere
\PassOptionsToPackage{unicode}{hyperref}
\PassOptionsToPackage{hyphens}{url}
%
\documentclass[
  10pt,
  a4paper,
]{article}
\usepackage{amsmath,amssymb}
\usepackage{iftex}
\ifPDFTeX
  \usepackage[T1]{fontenc}
  \usepackage[utf8]{inputenc}
  \usepackage{textcomp} % provide euro and other symbols
\else % if luatex or xetex
  \usepackage{unicode-math} % this also loads fontspec
  \defaultfontfeatures{Scale=MatchLowercase}
  \defaultfontfeatures[\rmfamily]{Ligatures=TeX,Scale=1}
\fi
\usepackage{lmodern}
\ifPDFTeX\else
  % xetex/luatex font selection
  \setmainfont[]{DejaVu Serif}
  \setmonofont[]{DejaVu Sans Mono}
  \setmathfont[]{Latin Modern Math}
\fi
% Use upquote if available, for straight quotes in verbatim environments
\IfFileExists{upquote.sty}{\usepackage{upquote}}{}
\IfFileExists{microtype.sty}{% use microtype if available
  \usepackage[]{microtype}
  \UseMicrotypeSet[protrusion]{basicmath} % disable protrusion for tt fonts
}{}
\makeatletter
\@ifundefined{KOMAClassName}{% if non-KOMA class
  \IfFileExists{parskip.sty}{%
    \usepackage{parskip}
  }{% else
    \setlength{\parindent}{0pt}
    \setlength{\parskip}{6pt plus 2pt minus 1pt}}
}{% if KOMA class
  \KOMAoptions{parskip=half}}
\makeatother
\usepackage{xcolor}
\usepackage[margin=20mm]{geometry}
\setlength{\emergencystretch}{3em} % prevent overfull lines
\providecommand{\tightlist}{%
  \setlength{\itemsep}{0pt}\setlength{\parskip}{0pt}}
\setcounter{secnumdepth}{-\maxdimen} % remove section numbering
\usepackage{amsmath,amssymb,amsthm,mathtools,graphicx,booktabs,longtable,needspace}
\usepackage{fvextra}
\DefineVerbatimEnvironment{Highlighting}{Verbatim}{commandchars=\\\{\},breaklines=true,breakanywhere=true,fontsize=\small}
\setlength{\emergencystretch}{3em}
\allowdisplaybreaks
\raggedbottom\sloppy
\widowpenalty=10000\clubpenalty=10000
\providecommand{\Tr}{\operatorname{Tr}}
\providecommand{\tr}{\operatorname{Tr}}
\ifLuaTeX
  \usepackage{selnolig}  % disable illegal ligatures
\fi
\usepackage{bookmark}
\IfFileExists{xurl.sty}{\usepackage{xurl}}{} % add URL line breaks if available
\urlstyle{same}
\hypersetup{
  pdftitle={Sharp two-coupling event region},
  pdfauthor={Kirtus Justus},
  hidelinks,
  pdfcreator={LaTeX via pandoc}}

\title{Sharp two-coupling event region}
\author{Kirtus Justus}
\date{}

\begin{document}
\maketitle

\subsection{Postulates}\label{postulates}

P1. A positive spectral mode at controlled relative coupling \texttt{c\textgreater{}0} produces an incident scale \texttt{U\textgreater{}0} with density

\texttt{f\_c(u)=4\ c\^{}2\ u\ exp(-2\ c\ u)}, \texttt{u\textgreater{}0}.

P2. For a registered event \texttt{B}, the detector has a measurable response \texttt{eta(u)=R(B\textbar{}u)} satisfying \texttt{0\textless{}=eta(u)\textless{}=1}, independent of \texttt{c}. A general noisy, nonlinear, lossy device is allowed. Record event counts \textbf{per prepared trial}, not conditional on successful detection.

P3. Coupling values and their ratios are separately calibrated, and the preparation law P1 is either established directly or generated by an unknown but \textbf{coupling-stable} mixture of positive modes whose eigenvalues all scale by the same factor \texttt{c}. Zero-mode-only trials, if present, must be handled with their missing-event mass explicitly.

P4. In the finite-shot theorem, trials within each coupling are independent Bernoulli experiments, and the detector event \texttt{B} is fixed before observing counts.

These are supplied physical/instrument assumptions, not consequences of the prephysical proposal.

\subsection{Theorem 1: sharp two-setting ROC region}\label{theorem-1-sharp-two-setting-roc-region}

Write \texttt{F\_c(t)=1-(1+2ct)exp(-2ct)} and consider couplings \texttt{1} and \texttt{b\textgreater{}1}. For \texttt{0\textless{}p\textless{}1} let \texttt{q\_p} be the unique positive root of \texttt{F\_1(q\_p)=p}. Define

\texttt{U\_b(p)=F\_b(q\_p)} and \texttt{L\_b(p)=1-F\_b(q\_(1-p))}, with continuous endpoint values \texttt{U\_b(0)=L\_b(0)=0}, \texttt{U\_b(1)=L\_b(1)=1}.

If \texttt{p\_c=Pr(B\textbar{}coupling\ c)} under P1-P2, then

\texttt{L\_b(p\_1)\ \textless{}=\ p\_b\ \textless{}=\ U\_b(p\_1)}.

\textbf{Sharpness:} for an isolated mode and any \texttt{p\_1}, the upper boundary is reached by the threshold detector \texttt{eta(u)=1\_\{u\textless{}=q\_p\}} and the lower boundary by \texttt{eta(u)=1\_\{u\textgreater{}=q\_(1-p)\}}. Every point between these boundaries at that fixed \texttt{p\_1} is attained by a convex combination of the two threshold kernels.

\textbf{Proof:} \texttt{f\_b(u)/f\_1(u)=b\^{}2\ exp{[}-2(b-1)u{]}}, strictly decreasing in \texttt{u}. Given a fixed \texttt{int\ f\_1\ eta=p\_1}, Neyman-Pearson's likelihood-ratio rearrangement maximizes \texttt{int\ f\_b\ eta} by placing detection weight on the smallest \texttt{u}, and minimizes it on the largest \texttt{u}. The threshold integrals are \texttt{F\_b(q\_p)} and \texttt{1-F\_b(q\_(1-p))}. Convex mixtures attain the interior. QED.

At \texttt{b=2}, \texttt{p\_1=1/2}, the admissible interval is exactly

\texttt{0.151832024412957731889...\ \textless{}=\ p\_2\ \textless{}=\ 0.848167975587042268111...}.

In particular, \texttt{(p\_1,p\_2)=(1/2,9/10)} is impossible for a stable detector acting on this known incident law.

\subsection{Theorem 2: robustness under an unknown stable spectral mixture}\label{theorem-2-robustness-under-an-unknown-stable-spectral-mixture}

Let a latent positive eigenvalue \texttt{Lambda} have \textbf{any fixed} probability distribution \texttt{pi}, independent of the controlled coupling \texttt{c}. Given \texttt{Lambda=lambda}, suppose

\texttt{f\_c(u\textbar{}lambda)=4(c\ lambda)\^{}2\ u\ exp{[}-2\ c\ lambda\ u{]}}.

The same envelope \texttt{L\_b(p\_1)\textless{}=p\_b\textless{}=U\_b(p\_1)} holds for an arbitrary stationary detector \texttt{eta(u)}, even though \texttt{pi} and the individual eigenvalues need not be known.

\textbf{Proof 1 (concavity):} conditional on \texttt{lambda}, make the substitution \texttt{v=lambda\ u}; Theorem 1 applies with the same functions \texttt{L\_b,U\_b}. Write the conditional event rates \texttt{p\_1(lambda),p\_b(lambda)}. The upper curve is strictly concave because \texttt{U\_b\textquotesingle{}(p)=b\^{}2\ exp{[}-2(b-1)q\_p{]}} is decreasing. The lower curve is convex because \texttt{L\_b(p)=1-U\_b(1-p)}. Jensen's inequality gives

\texttt{p\_b=E{[}p\_b(Lambda){]}\ \textless{}=\ E{[}U\_b(p\_1(Lambda)){]}\ \textless{}=\ U\_b(E{[}p\_1(Lambda){]})},

and analogously the lower bound. QED.

\textbf{Proof 2 (experiment randomization):} one can sample \texttt{V\textasciitilde{}f\_c}, draw \texttt{Lambda\textasciitilde{}pi} independently of \texttt{c}, and set \texttt{U=V/Lambda}. This is a coupling-independent Markov transformation of the original scale experiment. Passing \texttt{U} through the detector is a further randomization and therefore cannot exceed the original Neyman-Pearson ROC envelope. QED.

If the spectral weights change with coupling, or the detector changes with coupling, the conclusion does not follow. If an independently identified zero mode generates no registered positive-scale event, include that missing mass in the per-trial probabilities; the same universal envelope remains conservative by the concavity/convexity argument.

\subsection{Theorem 3: exact finite-dimensional detector-feasibility dual}\label{theorem-3-exact-finite-dimensional-detector-feasibility-dual}

For any known finite set of calibrated incident densities \texttt{f\_1,...,f\_m}, an event-probability vector \texttt{p=(p\_1,...,p\_m)} is attainable by one unchanged detector kernel \texttt{0\textless{}=eta(u)\textless{}=1} if and only if, for every real vector \texttt{a=(a\_1,...,a\_m)},

\texttt{sum\_j\ a\_j\ p\_j\ \textless{}=\ int\_0\^{}infty\ max(0,\ sum\_j\ a\_j\ f\_j(u))\ du}.

\textbf{Proof:} the map \texttt{eta\ -\textgreater{}\ (int\ f\_1\ eta,...,int\ f\_m\ eta)} sends the weak-* compact convex set \texttt{\{eta\ in\ L-infinity:\ 0\textless{}=eta\textless{}=1\}} into a compact convex subset of \texttt{R\^{}m}. Its support function in direction \texttt{a} is the stated integral, obtained by choosing \texttt{eta=1} wherever \texttt{sum\_j\ a\_j\ f\_j\textgreater{}0} and \texttt{eta=0} where it is negative. By finite-dimensional separating-hyperplane duality, the support inequalities characterize membership exactly. QED.

This theorem is a classical convex feasibility statement for \textbf{one event}, not a claim that checking a few selected directions certifies a common response kernel for an entire multibin detector.

At \texttt{m=2}, the continuum of support inequalities reduces to the sharp ROC boundaries above.

\subsection{Counterexample: all basic three-setting moment checks can pass while the two-setting envelope fails}\label{counterexample-all-basic-three-setting-moment-checks-can-pass-while-the-two-setting-envelope-fails}

For couplings \texttt{1,2,3}, consider the context-dependent binary detector with exact event probabilities

\texttt{p\_click=(1/2,9/10,19/20)} and \texttt{p\_miss=(1/2,1/10,1/20)}.

The normalized moment vectors under the proposed stationary Gamma model would be

\begin{itemize}
\tightlist
\item
  click: \texttt{(r0,r1,r2)=(1/2,9/40,19/180)};
\item
  miss: \texttt{(r0,r1,r2)=(1/2,1/40,1/180)}.
\end{itemize}

In \textbf{both} cases, \texttt{r0*r2-r1\^{}2=31/14400\textgreater{}0} and \texttt{r0\textgreater{}r1\textgreater{}r2\textgreater{}0}. Therefore the commonly used quadratic positivity, monotonicity, and conditional-variance conditions pass. Their total moments also add to \texttt{(1,1/4,1/9)}, the correct Gamma reference moments.

Nevertheless, \texttt{p\_click,2=0.9\textgreater{}U\_2(0.5)=0.8481679755...}. This known-reference bound rejects the stationary detector from \textbf{only the first two settings}.

A linear separating certificate is obtained at \texttt{q=q\_(1/2)} with slope \texttt{rho=4exp(-2q)=0.7467292354...}. The reference experiment obeys for every stationary detector

\texttt{p\_2-rho*p\_1\ \textless{}=\ U\_2(1/2)-rho/2\ =\ 0.4748033578...}.

The candidate violates this inequality by exactly \texttt{0.9-U\_2(1/2)=0.0518320244...}.

The example demonstrates that treating a small truncated Hausdorff moment test as complete detector certification is unsound when a specific input distribution and domination constraint are already known.

\subsection{Theorem 4: finite-trial falsification under unknown fixed spectral mixtures}\label{theorem-4-finite-trial-falsification-under-unknown-fixed-spectral-mixtures}

Choose couplings \texttt{1,2}, an event \texttt{B} selected in advance, and \texttt{M} independent prepared trials at each coupling. Let the empirical event fractions be \texttt{hat\ p\_1,hat\ p\_2}. For a target familywise error \texttt{alpha}, put

\texttt{eps=sqrt(log(4/alpha)/(2M))}.

Reject the stationary-detector null if

\texttt{hat\ p\_2-eps\ \textgreater{}\ U\_2(min(1,hat\ p\_1+eps))}

or if

\texttt{hat\ p\_2+eps\ \textless{}\ L\_2(max(0,hat\ p\_1-eps))}.

\textbf{Validity:} Hoeffding plus the two-proportion union bound shows that both true event probabilities lie within the displayed intervals with probability at least \texttt{1-alpha}. On this event, Theorem 2 prevents rejection. Hence the false-rejection probability is at most \texttt{alpha}, uniformly over all fixed spectral mixtures allowed by P3 and all unchanged detectors.

\textbf{Explicit power guarantee:} against the particular coupling-dependent alternative \texttt{(p\_1,p\_2)=(0.5,0.95)}, choose \texttt{alpha=0.05}, \texttt{M=2500} per setting. Then

\texttt{eps=0.029604143746...},

\texttt{0.95-2eps-U\_2(0.5+2eps)=0.002437333304...\textgreater{}0}.

On the same simultaneous Hoeffding concentration event, the upper-bound rejection is forced. This proves \textbf{type-I error at most 5\%, power at least 95\%, using 5,000 total independent preparations}, for the specified alternative. It is a conservative sufficient count, not an optimal trial count.

\end{document}
